\begin{table}[t]
\centering
\caption{\textbf{Two twin probes, and what closing the shape loophole does} (EQNET \texttt{poly8},
$K{=}500$, 3 seeds, chance $0.500$; tags \texttt{r72}, \texttt{r73}). Both twins have token and
operator multisets identical to the original, so inventory baselines are pinned at exactly $0.500$
by construction. The \emph{rotation} twin additionally changes depths, parent--child pairs and
subtree sizes (tier-3 cues), which is why the tree-local bag reaches $0.861$ there. The
\emph{swap} twin exchanges two siblings under a non-commutative operator, leaving the tree-local bag
exactly invariant, so shape is pinned too. Under the stricter probe the untrained encoder loses most
of its advantage while the trained encoder does not.}
\label{tab:twin}
\small
\begin{tabular}{l l rr}
\toprule
Baseline & Cue it can still see & Rotation & \textbf{Swap} \\
\midrule
Bag, full token (0 params)      & --- (pinned)            & $0.500$ & $0.500$ \\
Bag, tree-local (0 params)      & shape (rotation only)   & $0.861$ & $0.500$ \\
Bag, token n-gram (0 params)    & order                   & $0.559$ & $0.681$ \\
Untrained encoder (non-learned) & order, shape, recursion & $0.925$ & $0.788$ \\
\midrule
\textbf{Trained Tree-LSTM}      & learned                 & $\mathbf{1.000}$ & $\mathbf{0.994}$ \\
\midrule
\multicolumn{2}{l}{\emph{margin over the untrained architecture}} & $0.075$ & $\mathbf{0.206}$ \\
\multicolumn{2}{l}{\emph{classes probed}} & 395/500 & 348/500 \\
\bottomrule
\end{tabular}
\end{table}
